\section{Zadanie 2}
\autor{Hubert Jastrzębski}
(Spacer losowy z jedną barierą)
Startujemy doświadczenie z liczbą \(0\). Dopóki aktualna liczba jest mniejsza od \(n\),
dodajemy do niej \(1\) lub \(-1\) z równym prawdopodobieństwem.
Jakie jest prawdopodobieństwo, że dotrzemy do \(n\)?

\subsection*{Rozwiązanie}
Niech \(A_i\) to zdarzenie polegające na tym,
że doszliśmy do \(n\) przed dojściem do \(-i\).

Zauważmy, że \(\forall _{i\in\mathbb{N_+}} A_i \subset A_{i+1}\) \\
(ponieważ jeśli doszliśmy do \(n\) przed \(-i\), to na pewno doszliśmy do \(n\) przed \(-(i+1)\)).

Zatem z [Zestaw 1 Zadanie 2] mamy, że: 
\[
P\left(\bigcup_{i=1}^{\infty} A_i\right) = \lim_{i\to\infty} P(A_i)
\]

Chcemy teraz policzyć \(P(A_i)\).
Zauważmy, że możemy to przedstawić jako spacer losowy z dwoma barierami na pozycjach
\(0\) oraz \(n + i\), który zaczyna się na pozycji \(i\) 
(ponieważ przesuwamy wszystkie pozycje o \(+i\), więc jest to równoważne \(A_i\),
gdzie mamy bariery na \(-i\) oraz \(n\) i startujemy w \(0\)).
Natomiast z [Zestaw 1 Zadanie 10] wiemy, że coś takiego wynosi \(\frac{i}{n+i}\).

Tak więc:
\[
P\left(\bigcup_{i=1}^{\infty} A_i\right) = \lim_{i\to\infty} P(A_i) = \lim_{i\to\infty} \frac{i}{n+i} = 1
\]

Zostaje nam zastanowienie się czym jest \(\bigcup_{i=1}^{\infty} A_i\).
Jest to zbiór wszystkich trajektorii spaceru losowego, które w pewnym momencie osiągną poziom \(n\).
Jeśli więc prawdopodobieństwo na wystąpienie jakiegoś zdarzenia z tego zbioru wynosi \(1\),
to znaczy, że z prawdopodobieństwem \(1\) w pewnym momencie osiągniemy poziom \(n\).
\subsection*{Odpowiedź}
\[
\boxed{\text{1}}
\]
